% GENERATED FILE. Edit canonical lessons, then run npm run generate:books.
% canonical-source-hash: fnv1a64:024075e2125dc183
% canonical-lessons: FA-C486-hayajanzade, FA-C486-zoq-kardan, FA-C486-sarhal, FA-C486-eshq, FA-C486-motenaffer-budan

\chapter{Joy and Love}
\label{ch:fa-joy-and-love}
\hlchaptermodality{486}

\begin{chapteropening}
\textbf{By the end of this chapter:} \emph{I can talk about joy, love and dislike.}

Complete the last lesson of chapter 486: I can talk about joy, love and dislike.
\end{chapteropening}

\section[\texorpdfstring{\fa{هیجان} \fa{زده} (hayajânzade)}{هیجان زده (hayajânzade)}]{\fa{هیجان} \fa{زده} (hayajânzade) — excited}

\label{lesson:FA-C486-hayajanzade}



\begin{quote}
Before the new one: say the Persian for exasperated, then the Persian for to sigh.
\end{quote}

\subsection*{You'll want to know: \fa{هیجان} \fa{زده}}
\textbf{\fa{هیجان} \fa{زده}} (\emph{hayajânzade}) — \textquotedblleft{}excited\textquotedblright{}.

\begin{quote}
\fa{بچه} \fa{ها} \fa{برای} \fa{سفر} \fa{هیجان} \fa{زده} \fa{هستند}.
\end{quote}

\begin{quote}
\emph{bachchehâ barâ-ye safar hayajânzade hastand.}
\end{quote}

\begin{quote}
The children are excited about the trip.
\end{quote}

\emph{hayajân} 'excitement' + \emph{zade} 'struck' (from \emph{zadan}): struck by excitement.

\begin{cousinweb}
\emph{hayajânzade} is built in Persian on \emph{hayajân}, 'excitement', a loan usually traced to Arabic \emph{hayajān}, 'agitation', and \emph{zade}, 'struck', from the inherited verb \emph{zadan}. EtymDB has no entry for it.
\end{cousinweb}

\begin{grammarlens}[title={What you've built}]
One more word for telling, explaining and describing at B1.
\end{grammarlens}

\subsection*{Your turn}
\begin{itemize}
\raggedright
  \item \emph{Say it:} \emph{hayajânzade}
  \item \emph{Say it:} \emph{hayajânzade}, once more
  \item \emph{Say it:} a sentence of your own with \emph{hayajânzade}
  \item \emph{Recall:} say the Persian for exasperated, then the Persian for to sigh, then say \emph{hayajânzade} again
\end{itemize}

\begin{tcolorbox}[breakable,colback=teal!4,colframe=teal!35!black,title={Before you move on}]
What does \fa{هیجان} \fa{زده} mean? (\textquotedblleft{}Excited\textquotedblright{}.) Say it once more.
\end{tcolorbox}
\section[\texorpdfstring{\fa{ذوق} \fa{کردن} (zowq kardan)}{ذوق کردن (zowq kardan)}]{\fa{ذوق} \fa{کردن} (zowq kardan) — to be thrilled}

\label{lesson:FA-C486-zoq-kardan}



\begin{quote}
Before the new one: say the Persian for to sigh, then the Persian for excited.
\end{quote}

\subsection*{You'll want to know: \fa{ذوق} \fa{کردن}}
\textbf{\fa{ذوق} \fa{کردن}} (\emph{zowq kardan}) — \textquotedblleft{}to be thrilled\textquotedblright{}.

\begin{quote}
\fa{او} \fa{با} \fa{دیدن} \fa{هدیه} \fa{ذوق} \fa{کرد}.
\end{quote}

\begin{quote}
\emph{u bâ didan-e hediye zowq kard.}
\end{quote}

\begin{quote}
He was thrilled when he saw the present.
\end{quote}

\emph{zowq} 'delight; taste, flair' + \emph{kardan}. Warm and colloquial.

\begin{cousinweb}
\emph{zowq} is a loan, usually traced to Arabic \emph{ḏawq}, 'taste, relish'. EtymDB has no entry for it. \emph{kardan}, 'to do', is inherited Persian.
\end{cousinweb}

\begin{grammarlens}[title={What you've built}]
One more word for telling, explaining and describing at B1.
\end{grammarlens}

\subsection*{Your turn}
\begin{itemize}
\raggedright
  \item \emph{Say it:} \emph{zowq kardan}
  \item \emph{Say it:} \emph{zowq kardan}, once more
  \item \emph{Say it:} a sentence of your own with \emph{zowq kardan}
  \item \emph{Recall:} say the Persian for to sigh, then the Persian for excited, then say \emph{zowq kardan} again
\end{itemize}

\begin{tcolorbox}[breakable,colback=teal!4,colframe=teal!35!black,title={Before you move on}]
What does \fa{ذوق} \fa{کردن} mean? (\textquotedblleft{}To be thrilled\textquotedblright{}.) Say it once more.
\end{tcolorbox}
\section[\texorpdfstring{\fa{سرحال} (sarhâl)}{سرحال (sarhâl)}]{\fa{سرحال} (sarhâl) — in good spirits; lively}

\label{lesson:FA-C486-sarhal}



\begin{quote}
Before the new one: say the Persian for excited, then the Persian for to be thrilled.
\end{quote}

\subsection*{You'll want to know: \fa{سرحال}}
\textbf{\fa{سرحال}} (\emph{sarhâl}) — \textquotedblleft{}in good spirits; lively\textquotedblright{}.

\begin{quote}
\fa{امروز} \fa{صبح} \fa{خیلی} \fa{سرحال} \fa{هستم}.
\end{quote}

\begin{quote}
\emph{emruz sobh kheyli sarhâl hastam.}
\end{quote}

\begin{quote}
This morning I feel full of beans.
\end{quote}

\emph{sar} 'head' + \emph{hâl} 'state'. The opposite is \emph{bihâl}.

\begin{cousinweb}
\emph{sarhâl} is built in Persian on \emph{sar}, 'head', which EtymDB traces to Middle Persian (EtymDB), and \emph{hâl}, 'state', a loan from Arabic \emph{ḥāl}.
\end{cousinweb}

\begin{grammarlens}[title={What you've built}]
One more word for telling, explaining and describing at B1.
\end{grammarlens}

\subsection*{Your turn}
\begin{itemize}
\raggedright
  \item \emph{Say it:} \emph{sarhâl}
  \item \emph{Say it:} \emph{sarhâl}, once more
  \item \emph{Say it:} a sentence of your own with \emph{sarhâl}
  \item \emph{Recall:} say the Persian for excited, then the Persian for to be thrilled, then say \emph{sarhâl} again
\end{itemize}

\begin{tcolorbox}[breakable,colback=teal!4,colframe=teal!35!black,title={Before you move on}]
What does \fa{سرحال} mean? (\textquotedblleft{}In good spirits; lively\textquotedblright{}.) Say it once more.
\end{tcolorbox}
\section[\texorpdfstring{\fa{عشق} ('eshq)}{عشق ('eshq)}]{\fa{عشق} ('eshq) — love (passionate)}

\label{lesson:FA-C486-eshq}



\begin{quote}
Before the new one: say the Persian for to be thrilled, then the Persian for in good spirits.
\end{quote}

\subsection*{You'll want to know: \fa{عشق}}
\textbf{\fa{عشق}} (\emph{'eshq}) — \textquotedblleft{}love (passionate)\textquotedblright{}.

\begin{quote}
\fa{عشق} \fa{او} \fa{به} \fa{موسیقی} \fa{از} \fa{کودکی} \fa{شروع} \fa{شد}.
\end{quote}

\begin{quote}
\emph{'eshq-e u be musiqi az kudaki shoru' shod.}
\end{quote}

\begin{quote}
His love of music began in childhood.
\end{quote}

Strong love, for a person or a passion. For liking, Persian says \emph{dust dâshtan}.

\begin{cousinweb}
\emph{'eshq} was borrowed from Arabic \emph{ʿišq}, 'love, passion' (EtymDB).
\end{cousinweb}

\begin{grammarlens}[title={What you've built}]
One more word for telling, explaining and describing at B1.
\end{grammarlens}

\subsection*{Your turn}
\begin{itemize}
\raggedright
  \item \emph{Say it:} \emph{'eshq}
  \item \emph{Say it:} \emph{'eshq}, once more
  \item \emph{Say it:} a sentence of your own with \emph{'eshq}
  \item \emph{Recall:} say the Persian for to be thrilled, then the Persian for in good spirits, then say \emph{'eshq} again
\end{itemize}

\begin{tcolorbox}[breakable,colback=teal!4,colframe=teal!35!black,title={Before you move on}]
What does \fa{عشق} mean? (\textquotedblleft{}Love (passionate)\textquotedblright{}.) Say it once more.
\end{tcolorbox}
\section[\texorpdfstring{\fa{متنفر} \fa{بودن} (motenaffer budan)}{متنفر بودن (motenaffer budan)}]{\fa{متنفر} \fa{بودن} (motenaffer budan) — to hate}

\label{lesson:FA-C486-motenaffer-budan}



\begin{quote}
Before the new one: say the Persian for in good spirits, then the Persian for love.
\end{quote}

\subsection*{You'll want to know: \fa{متنفر} \fa{بودن}}
\textbf{\fa{متنفر} \fa{بودن}} (\emph{motenaffer budan}) — \textquotedblleft{}to hate\textquotedblright{}.

\begin{quote}
\fa{او} \fa{از} \fa{سرما} \fa{متنفر} \fa{است}.
\end{quote}

\begin{quote}
\emph{u az sarmâ motenaffer ast.}
\end{quote}

\begin{quote}
She hates the cold.
\end{quote}

\emph{motenaffer} 'averse' + \emph{budan} 'to be', used as \emph{az chizi motenaffer budan}. \emph{Motenaffer-am} means 'I hate it'.

\begin{cousinweb}
\emph{motenaffer} is a loan, usually traced to Arabic \emph{mutanaffir}, 'averse, recoiling'. EtymDB has no entry for it. \emph{budan}, 'to be', is inherited Persian.
\end{cousinweb}

\begin{grammarlens}[title={What you've built}]
One more word for telling, explaining and describing at B1.
\end{grammarlens}

\subsection*{Your turn}
\begin{itemize}
\raggedright
  \item \emph{Say it:} \emph{motenaffer budan}
  \item \emph{Say it:} \emph{motenaffer budan}, once more
  \item \emph{Say it:} a sentence of your own with \emph{motenaffer budan}
  \item \emph{Recall:} say the Persian for in good spirits, then the Persian for love, then say \emph{motenaffer budan} again
\end{itemize}

\begin{tcolorbox}[breakable,colback=teal!4,colframe=teal!35!black,title={Before you move on}]
What does \fa{متنفر} \fa{بودن} mean? (\textquotedblleft{}To hate\textquotedblright{}.) Say it once more.
\end{tcolorbox}
